\documentclass[11pt,a4paper]{article}
\usepackage[utf8]{inputenc}
\usepackage[T1]{fontenc}
\usepackage{lmodern,amsmath,amssymb,amsthm,mathtools,booktabs,hyperref,enumitem,microtype}
\usepackage{natbib}
\usepackage[margin=1in]{geometry}
\newtheorem{definition}{Definition}[section]
\newtheorem{theorem}[definition]{Theorem}
\newtheorem{proposition}[definition]{Proposition}
\newtheorem{corollary}[definition]{Corollary}
\theoremstyle{remark}\newtheorem*{remark}{Remark}
\title{The Law of Graphic Equalisation --- Mind and Self Grounding\\\large Frames, the Run and the Record, Free Will, Personal Identity, and Death}
\author{Matthew Parslow\\Independent Researcher}
\date{\textbf{Draft \textemdash\ draft-2026-08-27}\\[2pt]27 August 2026}
\usepackage{fancyhdr}
\pagestyle{fancy}
\fancyhf{}
\fancyfoot[C]{\thepage}
\fancyfoot[R]{\scriptsize\textsc{draft-2026-08-27}}
\renewcommand{\headrulewidth}{0pt}
\fancypagestyle{plain}{\fancyhf{}\fancyfoot[C]{\thepage}%
  \fancyfoot[R]{\scriptsize\textsc{draft-2026-08-27}}\renewcommand{\headrulewidth}{0pt}}
\begin{document}
\maketitle
\begin{center}
\fbox{\parbox{0.92\textwidth}{\small
\textbf{DRAFT --- draft-2026-08-27} (27 August 2026).
This document is under active revision and is circulated for comment, not as a
finished result. It is a companion to \emph{The Law of Graphic Equalisation ---
Theory} (referred to below as Graphic Equalisation) and to its physics companion
\emph{The Law of Graphic Equalisation --- Physics Grounding} (the physics
companion), and inherits their definitions, epistemic boundary, and
level-of-claims. Sections marked as falsifiability commitments are the intended
points of attack.}}
\end{center}
\vspace{1em}

\begin{abstract}
This companion to Graphic Equalisation develops the mind-and-self groundings of the
change-substrate. Its central move is that a frame is real graph data that
\emph{runs}, and that everything here follows from the run and its record. A rock
is a universe, so anything a rock does without is not compulsory for a universe.
Recording spatialises a succession into a static held axis, giving the run/record
axis whose poles are Conway (pure running, computationally irreducible) and a movie
(pure record), with real frames the machinima between; the past can be a record but
the future is uncomputed, so it is open as \emph{unrun}, not as indeterministic. Free
will is then the in-frame run --- a fully recorded, zero-randomness story
demonstrates it, its characters free in frame while the reader foreknows --- so
freedom needs no indeterminacy. Imagination is real: there is no non-actual frame,
and the only difference between the imagined and the perceived is frame-relative
grounding. Death is the self's rules ceasing to be re-runnable in the same coherent
self; personal identity is internal coherence (with an instance a supergraph node),
which places gradual replacement, teleporter, backup, and fission on one axis; and
death is relative to a chosen coherence-boundary, with many real deaths and none
absolute. The claims throughout are information-level and structural.
\end{abstract}

\section{Scope}\label{sec:scope}
The method is to treat philosophy of mind as applied physics: once a frame is graph
data that runs, free will, identity, and death are read off the run and its record,
not argued as separate doctrines. This paper is a grounding, not a psychology and not
a metaphysics asserted from within. It inherits Graphic Equalisation's position that the claims are about
information --- distinction, frame, admissibility, realisation, and now the
\emph{run} of a frame and its \emph{record} --- of which mind is one embodiment.
The strong claims are structural: that free will is the in-frame run and needs no
indeterminacy; that imagination is real and differs from perception only in
grounding; that personal identity is internal coherence (Graphic Equalisation's SELF
identity), distinct from the supergraph node that instances it (its SUPERGRAPH
identity); and that death is the loss of re-runnability in
the same coherent self, relative to a chosen coherence-boundary. The weak part is
anything operational: no decision procedure for a given choice, no test for when
re-runnability or coherence has in fact been lost, no threshold for how gradual a
change stays ``the same'' self, and no account of the subjective, felt, or
phenomenal character of any of it --- consciousness and qualia are out of scope.
Each grounding is a \emph{promoted} constraint exposed to falsification, in the
manner of Graphic Equalisation's collapse bridge, not a result asserted from within.

\section{The minimal universe}\label{sec:minimal}
A \emph{rock}, in the sense used throughout, is an inert frame: a graph whose rules
are present but not firing, so it has no succession and no time of its own. Things
happen to it from its supergraph, but in its own frame it merely is (the physics
companion develops this as why a rock has no time). Everything here is a graph, so a
rock and a universe are one kind of thing, differing in scale, not in category
(Graphic Equalisation's self-similarity: the runner is graph-that-runs at every
scale). A rock \emph{is} a universe --- the minimal one --- and that makes it an
existence proof about universes in general.

\begin{proposition}[What the rock shows is not compulsory]\label{prop:minimal-universe}
Anything not compulsory for a rock is not compulsory for a universe. The rock does
without time (it is inert, no succession), without cycles or closure, without
inference or observation, without history or records, without self-originated
change, and without a running arrow --- so none of these is compulsory for a
universe.
\end{proposition}

The inference is one-directional, and only that direction is claimed. A single inert example
bounds the compulsory set from \emph{above} --- it shows what can be dispensed with --- and
cannot establish that what the rock happens to retain is necessary. So the rock is a
non-necessity argument, not a derivation of a core.

What is compulsory is inherited from Graphic Equalisation, which keeps almost nothing: a
container relation across which rules run. The rock is consistent with that and does not
weaken it, because a rock's rules are inert \emph{in its own frame} while its container's are
not --- what happens to the rock (weathering, being moved) is its supergraph's rules running
across exactly that container relation. An inert frame is therefore a subgraph whose own rules
do not fire, not a counterexample to rules running; and it is in that sense, rather than as a
free-standing cosmos, that it is called a minimal universe here.

This is the general form of ``time is optional,'' and it licenses the
contingent-not-necessary verdicts of the physics companion. Wherever a shape there
was called contingent rule-content rather than necessity --- the number of spatial
dimensions, the gauge group, the dimensionless constants, the arrow, the value of
$c$ --- the rock is the reason: the minimal universe does without it, so it cannot
be compulsory. Physics ordinarily takes time, causality, and the second law's arrow
as features a universe must have; the rock shows they are features an \emph{active}
universe has.

Honest boundary: the argument is as strong as the self-similarity it rests on
(rock and universe as one kind, differing by scale). It does not claim a universe
\emph{is} inert like a rock --- an active universe has time, change, an arrow. It
claims only that a universe \emph{need not} have what the rock does without: those
features are contingent on activity, not compulsory to being a universe at all.

\section{Running and recording: Conway, a movie, and the machinima between}\label{sec:runrecord}
Recording spatialises time. A succession held as a record --- an MRI of a rock, a
film strip, a worldline drawn on paper --- lays the temporal order out as a static
structure, a $+$-like axis (the block's extra dimension). This reads the metric
signature: the $-$ axis is the one being \emph{lived}, the frame's own succession
run by its cycles, open at the frontier; the $+$ axes are \emph{held} --- spatial
positions, and any recorded time. The sign tracks lived-versus-held, not a
property intrinsic to the axis; the block feels dead because a recorded succession
is not run.

The extremes of the run-versus-record axis are Conway and a movie. Conway's Game
of Life is pure running: a rule and a seed, everything computed live, nothing
recorded, and no way to know the outcome but to run it (computational
irreducibility). A movie is pure record: every frame laid out, played back,
nothing computed --- the block universe, known by reading rather than running.
Between them is scripted machinima --- a running engine following a predetermined
script, part computed live, part laid down in advance --- and that is where real
frames sit: rules fire (Conway-like) while history is held (movie-like). It is the
same axis as generation-versus-record: store the generator and run it (Conway,
small but irreducible) or store the frames and read them (a movie, large but
immediate), with the world's own compressibility deciding which is cheaper.

\begin{proposition}[The past can be a movie; the future is Conway]\label{prop:run-record}
A block is only ever of the past. What has been run can be recorded and laid out
--- a movie of the succession so far, readable without re-running. What has not
been run cannot be recorded, because there is nothing yet to record: the future is
determinable by the rules but \emph{uncomputed}, Conway-style, so it is not laid out
and cannot be read ahead --- you must run it. The open frontier of Axiom~$-1$ is
therefore not indeterminism but unrun-ness: the future is open because it is
uncomputed, not because nothing fixes it, and you cannot record what has not run.
\end{proposition}

\section{Free will and the three frames}\label{sec:freewill}
\begin{proposition}[Free will is the in-frame run; a story demonstrates it]\label{prop:free-will}
Free will is a frame running its own decision. In a frame's own run the decision is
live and unrun-to-it --- the frame must run its deliberation to reach the outcome ---
and that running \emph{is} the choosing: the outcome is fixed by the rules yet
knowable only by running, so the frame experiences computing it as choosing it. This
needs no indeterminacy, and a story is the perfect demonstration. A book or a movie
is a fully recorded succession --- every word or frame fixed, zero randomness,
readable ahead by the reader --- and yet its characters have complete free will:
they deliberate, choose, and are responsible, in frame. Both hold at once because
they are two frames. The character's freedom is its in-frame run; the reader's
foreknowledge is a superframe holding the finished record (a movie), and holding the
record never reaches into the character's run to remove the freedom. So foreknowledge
is not unfreedom --- not even complete foreknowledge, the reader who knows the ending
or the author who wrote it --- because knowing the recorded run and being the live run
are different frames. Randomness would only subtract: a random run is not more free,
it is less the character's own. And responsibility survives for the reason the freedom
does --- the run is the character's own (self-change is the only change; the rules
running are the character) --- so a determined, recorded, even foreknown act is still
that character's act.
\end{proposition}

\begin{remark}[Author, character, reader: unrun-ness is absolute]\label{rem:three-frames}
A story stacks three frames on one succession. The author \emph{runs} it --- writing is
Conway, generated live, and the author must write to know the ending (authorship is
computationally irreducible; the commonplace that ``the characters took on lives of
their own'' is the author unable to read ahead their own creation). The character
\emph{lives} it --- the in-frame run of Proposition~\ref{prop:free-will}. The reader
\emph{replays} it --- the finished record, a movie. One succession, three
frame-readings: the run/record axis is a fact about frame, not about the system. And
the author settles the objection to Proposition~\ref{prop:run-record} that a high
enough superframe might read the future ahead --- even the creator gets no shortcut.
The author is Conway to their own book while writing it; the movie exists only
\emph{after} the run, for everyone including the author. Foreknowledge is always
post-run record-holding, never pre-run read-ahead. Unrun-ness is absolute: no frame,
creator included, reads an unrun future; every frame either runs it or replays an
already-run record.
\end{remark}

\section{Imagination, grounding, and reality}\label{sec:grounding}
\begin{proposition}[There is no non-actual frame; imagination is real]\label{prop:imagination-real}
Imagination is not the running of a non-actual frame --- there is no such thing. Every
frame that runs is real graph data really running, so actuality is not a property some
frames have and others lack; every frame is actual. An imagined dragon, a planned move,
a character in a book is a real frame, really run (real computation), not a pretend one.
The only difference between an imagined frame and a perceived one is frame-relative
grounding: what the frame is anchored to, relative to another frame. A perception is
grounded in the observation frame (sensory-anchored); an imagination is grounded
internally, in generation or in a held or communicated record. Both are real; they
differ in what they are grounded in, not in whether they are. So ``is it real?'' is
malformed as a question about actuality --- everything that runs is --- and the
well-formed question is ``what is it grounded in?'': observation, memory, generation, or
a communicated record. A shared fiction is a real
frame grounded in a communicated record (the book of Remark~\ref{rem:three-frames}); a
plan is a real frame grounded in projection. This is why imagination does work ---
planning, design, and counterfactual reasoning are real frames really run, not pretence
--- and why fiction moves: the character's frame runs, really, in the reader. Death, in
the same terms, is not a frame going non-actual (nothing does) but a run that stops,
leaving frames grounded only in a record and re-runnable by others; reading and
remembering are that re-running --- a record-grounded frame run again into a live
grounding.
\end{proposition}

\begin{proposition}[Grounding has no bedrock; observation is supergraph-contact]\label{prop:grounding}
Grounding is frame-relative all the way down: a frame grounded in a frame grounded in a
frame, with no absolute, self-grounding frame beneath them all. There is no bedrock
``actual'' layer --- the ground is the coherent running structure, frame to frame, not a
base. Yet one grounding is privileged, relationally and not absolutely: \emph{observation}
is a frame's contact with its supergraph, the boundary where the container deforms it with
changes it does not author (Graphic Equalisation's dissolved pointer --- you change
yourself, the supergraph deforms you --- with perception the deformation side). Observation
is special because it is \emph{yours}, your container-boundary, not because it is ``the
real.'' It follows that reality-testing is not actuality-checking --- there is no
non-actual frame to catch (Proposition~\ref{prop:imagination-real}) --- but
\emph{supergraph-pushback}: does the container deform you as if the frame were grounded
there? You walk into the chair (pushback: observation-grounded) and not into the dragon
(no pushback: imagination-grounded), both real frames, only one with the container pushing
back. Sanity is then a grounding-tag problem, not a reality one: keep observation-grounded
frames tagged as observation-grounded and imagination-grounded ones as
imagination-grounded. A hallucination or delusion is a grounding-tag error --- a real,
internally-grounded frame mistagged as observation-grounded, so the self expects a pushback
that will not come --- not a belief in something non-actual, since nothing is. The position
is thus anti-foundationalist (no absolute ground) without being flatly coherentist
(observation is privileged, as your supergraph-contact).
\end{proposition}

\section{Self-models: incomplete by construction, calibrated by choice}\label{sec:selfmodel}
A frame does not read its own trajectory. It \emph{carries} it --- the standing structure is what the frame is, not something it holds at a distance --- so to treat its own course as an object at all, a frame must build a model of itself. Three features of that model follow from the substrate rather than from any fact about minds.

It is \emph{incomplete}, because it is built under a finite aperture, and what falls outside is \textsc{null} rather than zero: an absent distinction does not present itself as a blank to be filled, it does not present at all. It is \emph{revisable only at a price}, because altering standing structure is the paid operation while retaining it is free (Narrative momentum), so coherence-preserving readings are cheap and re-writes are not. And it is \emph{indirect about its own subject}, since a frame is readable from outside in a way it is not from inside --- which is why another frame may model one's trajectory more accurately than one does oneself, for reasons that are structural and require no appeal to self-deception.

\begin{proposition}[Incompleteness is structural; awareness of it is not precluded]\label{prop:selfmodel}
Every self-model is incomplete, and no frame can escape this by effort or honesty. But a forced aperture is not an invisible one. A boundary leaves a trace on the inside --- residual: what fails to resolve, what arrives unexplained, what a rule set repeatedly cannot absorb --- and from that trace a frame can infer the \emph{fact} and often the \emph{shape} of its own limit, exactly as it infers its container from imposed change. The content beyond the aperture stays unavailable; that there is content, and roughly where the edge runs, need not. So the operative distinction is not between complete and incomplete self-models, all being incomplete, but between \emph{calibrated} and \emph{uncalibrated} ones.
\end{proposition}

The asymmetry, though, is one of \emph{access mode}, not of accuracy, and the two must not be run together. Mode is only one input to how good a model is; machinery and sustained investment are another, and they can dominate it. An outside frame has the better access and almost always the lower investment --- it is not working at the problem --- while a frame modelling itself has the worse access but continuity, sustained effort, and data no outsider obtains, its own residual among them. So the familiar case in which others judge one more accurately than one judges oneself is the \emph{common} case rather than the structural outcome, and it holds because most self-models are uncalibrated, not because the arrangement forces it. A sufficiently skilled frame can hold a better model of itself than any outside frame holds of it --- and, since building the model is itself the objectification, that skill converts part of what was merely carried into something read: the vantage a frame does not natively have, it can construct. The same machinery improves its models of others, which is why the two capabilities travel together rather than trading off.

This bears directly on how such failures should be read. A frame with no meta-machinery fills gaps it cannot see as gaps, and reports the filling as knowledge; a frame that models its own boundary reports the gap instead. The difference is a capability, acquirable and losable, not a virtue --- and, like everything else here, it is optional: a frame may hold a model of its own limits, or may not.

Honest boundary: this claims a shape --- that incompleteness is the default condition of any frame modelling itself, that it is structural rather than pathological, and that calibration against it is possible without removing it. It supplies no mechanism for any particular psychological phenomenon, no threshold at which a model counts as calibrated, and nothing about the felt character of self-knowledge, which remains outside this paper's scope.

\section{Death and personal identity}\label{sec:death}
\begin{proposition}[Death is when the self's rules cease to be re-runnable in the same coherent self]\label{prop:death}
Death is, critically, the rules within the self ceasing to be re-runnable in the same
coherent self entity. Not every stop is death: a pause --- sleep, anaesthesia, a
recoverable coma --- is a stop whose rules remain re-runnable in the same coherent
self, and the self resumes as itself. Death is the stop that does not resume: the
coherence that made these rules one self is gone, so they cannot be re-run as that
self, even if the graph-data they were persists scattered as a record. This is the
identity criterion --- the self persists exactly while its rules stay re-runnable as
the same coherent entity --- and it is why a pause is survivable and death is not; the
self-graph goes inert for good, a record where it was a run (the rock of
Section~\ref{sec:minimal}, now grounded only in what remains of it).

What ends is only the self's own run. The supergraph runs on, and two things the self
did persist in it. Its changes persist: rule-actions are pointers, work done once, so
every write the agent made to the supergraph is the supergraph's now, and death does not
undo it --- though under no-copy it persists only as surviving present structure, until
overwritten, not as an indelible record. And its copies persist: every frame the agent communicated,
taught, wrote, or built runs on in other frames, grounded there in the record it left
(the reader re-running the author of Remark~\ref{rem:three-frames}). But those copies
are partial fragments, not a full coherence-snapshot: they run as fragments in other
coherent selves, not as this coherent self re-run, so their persistence is not the
self's continuation (a full snapshot is the different case of Proposition~\ref{prop:identity},
where the self does re-run, forked to a new node). The
self can no longer self-change, and self-change was the only change it had; the changes
it already made and the copies already running elsewhere are no longer its to stop.
They keep running without it.
\end{proposition}

\begin{proposition}[Identity is internal coherence; instances are supergraph nodes]\label{prop:identity}
Two things must be kept apart: the self's \emph{internal coherence} --- its rules and
state, which is what Graphic Equalisation calls SELF identity --- and its
\emph{supergraph node} --- the instance of it running at a position in the container,
which is what that paper's SUPERGRAPH identity tracks. The self is the internal coherence; a node is one instance of it. This
resolves the standard puzzles. Gradual replacement (cell turnover, the Ship of Theseus)
is survival: internal coherence is continuous and there is one node throughout --- the
self is the coherent running, not the substrate. A backup restored is still the same
internal self --- a full coherence-snapshot re-run reinstates that exact coherence ---
but now at a \emph{second} supergraph node, and from the restore point the two diverge,
each accruing its own changes under its own grounding. That is a fission: one internal
self, two supergraph nodes, diverging, both the self at the fork and neither uniquely
the node after it --- which is Graphic Equalisation's own statement that a copy can
preserve SELF identity while receiving a new SUPERGRAPH identity, read as a result
about persons. A teleporter is the same with one branch pruned --- the full snapshot
reinstates the internal self at a new node, the original node stops --- so the internal
self survives, relocated. What matters about a record is whether it is \emph{full} or
\emph{partial}: a full coherence-snapshot re-runs into the same internal self (forked to
a new node), while a partial record --- a corpse, a memory, a written portrayal ---
cannot reconstitute the coherence, so re-running it starts a new coherent self, not the
original. The internal self therefore persists exactly as long as a re-runnable full
coherence-snapshot exists somewhere; the death of a node is one instance stopping, and
the death of the self is that re-runnability being gone everywhere.
\end{proposition}

\begin{remark}[Death is relative to the coherency of the self]\label{rem:death-relative}
It follows that death is relative to the coherency of the self. The criterion applies
within a coherent self, but which coherence counts as ``the same self'' is a
frame-relative choice, so death is indexed to that choice, not absolute. Drawn narrowly
--- this instant's exact coherent state --- the self dies each moment, replaced by the
next; drawn broadly --- the coherent person tolerant of change across a life --- it dies
once, at the irreversible break. A deep transformation (``I am not the person I was'')
is a death relative to the narrow coherence and a survival relative to the broad one,
both true at once. There is no absolute self, so there is no absolute death --- only a
death per coherence-boundary. This makes death relative, not unreal: each boundary's
death is real for that boundary --- the broad life-self's coherence does break
irreversibly, and that death is as real as any --- so relativity means many real deaths
at many boundaries, never none.
\end{remark}

\section{Positioning and prior art}\label{sec:prior-art}
The positions here are not new as positions; what is new is that each falls out of one
substrate mechanism --- the run, and its record --- rather than being posited on its
own. On free will this is a compatibilism, determinism and freedom coexisting in the
tradition of \citet{hume1748} and \citet{dennett1984}, but with a specific mechanism for
the freedom: the in-frame run and its computational irreducibility \citep{wolfram2002},
so that the future is open as \emph{unrun}, not as undetermined, and a fully recorded,
zero-randomness story still exhibits free will. On personal identity the continuity
criterion is Lockean in spirit \citep{locke1690} and closest to \citet{parfit1984} ---
identity as psychological or functional continuity, fission that branches, and what
matters being the continuity rather than a bare identity fact --- with the addition of
the split between internal coherence (the self) and supergraph node (the instance),
which places teleporter, backup, and fission on one axis. On time the past-as-record /
future-as-run asymmetry engages the block-universe debate --- to which \citet{putnam1967} argues
the block conclusion directly from special relativity, while \citet{mctaggart1908} supplies the
A-series/B-series apparatus in which the question is posed rather than the block thesis itself:
a block of the past is an ordinary record, while a completed block with a fixed future
is not, because the future is uncomputed. And the no-absolute-self reading of death ---
a death per coherence-boundary, the narrow self dying each moment --- recovers a
momentariness familiar from other traditions, here as a consequence of frame-relativity
rather than a doctrine. The distinctive claim throughout is the derivation, not the
destination.

One position here is older than the rest, and has been made domain by domain: the
refusal to treat any frame --- human frames included --- as privileged
(Section~\ref{sec:no-privilege}). In cosmology it is the
Copernican or mediocrity principle --- that Earth, and our position generally, is not privileged
--- named for the displacement of Earth in \citet{copernicus1543}, though the principle itself was
drawn out by successors rather than stated there. In physics it is the principle of relativity --- no privileged
reference frame \citep{einstein1905} --- and, for the observer specifically, relational
quantum mechanics, in which states are relative to any physical system and no observer is
special \citep{rovelli1996}. On free will it is Spinoza's: humans are not exempt from
nature's necessity, and feel free only through ignorance of their causes
\citep{spinoza1677}. On the self it is a long anti-substantial line --- Hume's bundle in
place of a persisting ego \citep{hume1739}, the Buddhist \emph{anatta} that \citet{parfit1984} drew on,
Dennett's self as a centre of narrative gravity rather than an inner witness
\citep{dennett1991}, Metzinger's phenomenal self-model with no one behind it
\citep{metzinger2003}, and Hofstadter's substrate-independent self-as-pattern
\citep{hofstadter2007}. On knowledge it is the rejection of a privileged foundation ---
Sellars against the Given \citep{sellars1956}, Rorty against the mind as a mirror
\citep{rorty1979}, Putnam against the God's-Eye View \citep{putnam1981} --- the same move
this paper's grounding-without-bedrock makes. Biology had already removed the species
privilege \citep{darwin1871}. And the standing opposition is equally identifiable: the
claim that the subjective point of view is irreducible, from \citet{nagel1974} and developed in
\citet{nagel1986}, is exactly the assumption these dissolutions decline --- though Nagel himself
is careful that this shows only that we lack a conception of how physicalism could be true, not
that it is false, so the privilege is what the present account reads in that irreducibility
rather than a thesis he asserts.

The move made here is the same one --- one more refusal of privilege --- and the lineage
above is best read as an asymptotic approach: each source removes the privilege from one
more domain, the sequence tending toward a limit of no privilege anywhere without any single
step reaching it, and, like any frontier, never completed --- there is always one more
privilege still unnoticed. Against that, the contribution here is both tiny and different in
kind. Tiny, because it is not a grand new principle: it drops the same assumption the others
drop. Different in kind, because it is not one more step along the asymptote --- not another
domain de-privileged --- but the substrate on which the refusal becomes a single uniform
operation rather than a domain-by-domain series. That is orthogonal to the sequence, not
further along it: naming the move the whole asymptote was making, and giving it one ground
to run on. The honest size of the claim is exactly that --- small, and of a different type
from the components it unifies.

\section{Conclusion: the one assumption removed}\label{sec:no-privilege}
The simplicity of every dissolution above comes from a single refusal: that humans ---
or any frame --- are privileged. It is the most pervasive assumption in these problems
and the least noticed, and each hard problem was made hard by smuggling it in. The
measurement problem privileges a human (or conscious) observer as the special coupling
that collapses; the flow-versus-block dispute privileges human-felt passage as the real
time; the free-will problem privileges human choice as needing a metaphysical freedom no
other system has; the identity and death problems privilege the human self as a persisting
soul rather than a coherent run; the reality-versus-fiction and foundation problems
privilege human perception as access to ``the real.'' Remove the privilege --- treat a
human frame by the same frame-relative rules as a rock, a story character, or a universe
--- and each collapses into the general machinery already built: any coupling measures,
time is any active frame's cycle-count, free will is any frame's in-frame run, the self is
internal coherence, and observation is one frame's supergraph-contact rather than a window
on the real. The physics companion's dissolutions (the special observer, the initial
state) are the same refusal in the other domain.

The assumption persists because it is uncomfortable to drop. Unprivileged means not the
centre, mortal as the loss of a coherence, free only as any frame is free, and never in
contact with ``the real'' except as one's own container-boundary. The discomfort is why
the assumption is both pervasive and invisible: it is motivated, held rather than observed.
So the framework is simple not because its answers are clever but because the difficulty
was manufactured by the privilege, and removing it removes the difficulty. This is
Graphic Equalisation's Axiom~$-1$ applied to observers: no frame is the god's-eye, a human
frame makes no privileged difference, and ``you are here'' is a location, not a throne.

Honest boundary: this is a methodological claim about \emph{why} the dissolutions are
simple, not a further mechanism, and it is not the claim that human frames are nothing.
Human frames are distinctive --- rich aperture, language, self-models --- but distinctive
is not privileged: being one remarkable frame among many is not being the exception to the
rules the rest of them run under. No-privilege removes the exemption, not the richness.

\section{Falsifiability commitments}\label{sec:falsifiability}
\begin{itemize}[leftmargin=1.4em]
\item \textbf{Compulsory versus contingent.} Fails if some feature the rock
demonstrably lacks (time, history, self-change, a running arrow) is nonetheless shown
compulsory for any graph to exist --- if the minimal universe (an inert graph) is
incoherent as a graph. The argument is only as strong as the self-similarity it rests
on (rock and universe one kind, differing by scale).
\item \textbf{The block is always of the past.} Fails if a completed block with a fixed
future is coherent under Axiom~$-1$, or if a block of the succession-so-far (an MRI of a
rock) is forbidden. The claim predicts records only ever trail the run, never lead it.
\item \textbf{The future is unrun, not unfixed.} Fails if a determinable future can be
read ahead without running it --- a shortcut past computational irreducibility for a
genuinely irreducible rule --- or if a block can contain a not-yet-run future.
\item \textbf{Free will is the in-frame run, not indeterminacy.} Fails if a fully
recorded, zero-randomness system (a story) cannot exhibit the form of free will ---
deliberation, choice, responsibility --- in its characters' frames, or if a
superframe's complete foreknowledge demonstrably removes a subframe's in-frame
freedom. The claim predicts foreknowledge and freedom coexist across frames.
\item \textbf{Actuality is running, not a frame-property.} Fails if some running
frame is demonstrably non-actual (real graph data that runs yet is not real), or if
the imagined/perceived distinction requires an actuality property rather than
frame-relative grounding. The claim predicts imagination does real computational work
(planning, design) as real runs, not pretence.
\item \textbf{Grounding has no bedrock; reality-testing is supergraph-pushback.} Fails
if there is an absolute, self-grounding frame beneath all others, or if reality-testing
can distinguish a real frame from a non-actual one by anything other than whether the
container (supergraph) pushes back. The claim predicts hallucination and delusion are
grounding-tag errors, not false beliefs about actuality.
\item \textbf{Death is loss of re-runnability in the same coherent self.} Fails if a
pause that preserves re-runnability as the same coherent self (sleep, a recoverable
coma) counts as death, or a stop that loses it counts as survival --- the criterion is
re-runnability in the same coherent entity, not the mere stopping. Also fails if a
self's writes to the supergraph vanish when its run stops (work-done-once violated), or
if frames it copied into other frames cease when it does. The claim predicts
consequences and copies persist, that copies are not the self resuming, and that
nothing re-runs as the same coherent self.
\item \textbf{Self-models are incomplete but calibratable.} Fails if a frame can hold a complete model of itself, or --- in the other direction --- if a forced aperture is necessarily invisible from the inside, so that no frame can represent the fact or shape of its own limit. The claim predicts that residual is the trace from which a boundary is inferred, that calibration is an acquirable capability rather than a fixed property, and that a sufficiently skilled frame can model itself better than an outside frame does --- so it fails too if outside access always beats inside investment, making the common case a structural necessity.
\item \textbf{Identity is internal coherence; death is relative to the
coherence-boundary.} Fails if the self is the substrate parts rather than the internal
coherence (gradual replacement must survive; a full coherence-snapshot restored must be
the same internal self forked to a second node, diverging; a partial record re-run must
be a new self; a teleporter must relocate the same internal self). And fails if there is
a single absolute self whose end is the one death independent of any coherence-choice, or
if no coherence-boundary's death is real --- the claim predicts many real deaths at many
boundaries and levels (node versus self), none absolute.
\item \textbf{No frame is privileged.} Fails if any dissolution here secretly requires a
human frame to be privileged --- a special observer, a special time, a special freedom, a
persisting soul, a window on the real --- rather than following from the frame-relative
machinery applied uniformly to every frame. The claim is that removing the privilege, not
any added mechanism, is what dissolves each problem.
\item \textbf{Scope.} The paper claims shapes only; it is falsified as over-claiming if
any section is read as a decision procedure, an operational test of when
re-runnability or coherence has been lost, a threshold for how gradual a change stays
the same self, or a theory of consciousness or qualia --- none is offered.
\end{itemize}

\begin{thebibliography}{99}
\bibitem[Copernicus(1543)]{copernicus1543}
Nicolaus Copernicus.
\newblock \emph{De revolutionibus orbium coelestium, libri VI}.
\newblock Nuremberg: Johannes Petreius, 1543.

\bibitem[Darwin(1871)]{darwin1871}
Charles Darwin.
\newblock \emph{The Descent of Man, and Selection in Relation to Sex}, 2 volumes.
\newblock John Murray, London, 1871.

\bibitem[Dennett(1984)]{dennett1984}
Daniel C. Dennett.
\newblock \emph{Elbow Room: The Varieties of Free Will Worth Wanting}.
\newblock MIT Press, Cambridge, MA, 1984.

\bibitem[Dennett(1991)]{dennett1991}
Daniel C. Dennett.
\newblock \emph{Consciousness Explained}.
\newblock Little, Brown, Boston, 1991.

\bibitem[Einstein(1905)]{einstein1905}
Albert Einstein.
\newblock Zur Elektrodynamik bewegter K\"orper.
\newblock \emph{Annalen der Physik} (ser.~4), 17:891--921, 1905. (Whole-series numbering: 322(10).)

\bibitem[Hofstadter(2007)]{hofstadter2007}
Douglas R. Hofstadter.
\newblock \emph{I Am a Strange Loop}.
\newblock Basic Books, New York, 2007.

\bibitem[Hume(1739--40)]{hume1739}
David Hume.
\newblock \emph{A Treatise of Human Nature}, Book~I, Part~IV, \S6 (``Of Personal Identity'').
\newblock London: John Noon, 1739--40.

\bibitem[Hume(1748)]{hume1748}
David Hume.
\newblock \emph{An Enquiry Concerning Human Understanding}, \S VIII (``Of Liberty and Necessity'').
\newblock London: printed for A. Millar, 1748. First published 1748 as \emph{Philosophical Essays Concerning Human Understanding}; retitled 1758.

\bibitem[Locke(1689/1694)]{locke1690}
John Locke.
\newblock \emph{An Essay Concerning Human Understanding}, Book~II, Chapter~xxvii (``Of Identity and Diversity'').
\newblock Edited by Peter H. Nidditch, Clarendon Edition of the Works of John Locke. Clarendon Press, Oxford, 1975. Originally 1689/1694; the identity chapter was added in the second edition (1694) and is not in the first.

\bibitem[McTaggart(1908)]{mctaggart1908}
J.~M.~E. McTaggart.
\newblock The unreality of time.
\newblock \emph{Mind}, 17(68):457--474, 1908.

\bibitem[Metzinger(2003)]{metzinger2003}
Thomas Metzinger.
\newblock \emph{Being No One: The Self-Model Theory of Subjectivity}.
\newblock MIT Press, Cambridge, MA, 2003.

\bibitem[Nagel(1974)]{nagel1974}
Thomas Nagel.
\newblock What is it like to be a bat?
\newblock \emph{The Philosophical Review}, 83(4):435--450, 1974.

\bibitem[Nagel(1986)]{nagel1986}
Thomas Nagel.
\newblock \emph{The View from Nowhere}.
\newblock Oxford University Press, New York, 1986.

\bibitem[Parfit(1984)]{parfit1984}
Derek Parfit.
\newblock \emph{Reasons and Persons}.
\newblock Clarendon Press, Oxford, 1984.

\bibitem[Putnam(1967)]{putnam1967}
Hilary Putnam.
\newblock Time and physical geometry.
\newblock \emph{The Journal of Philosophy}, 64(8):240--247, 1967.

\bibitem[Putnam(1981)]{putnam1981}
Hilary Putnam.
\newblock \emph{Reason, Truth and History}.
\newblock Cambridge University Press, Cambridge, 1981.

\bibitem[Rorty(1979)]{rorty1979}
Richard Rorty.
\newblock \emph{Philosophy and the Mirror of Nature}.
\newblock Princeton University Press, Princeton, NJ, 1979.

\bibitem[Rovelli(1996)]{rovelli1996}
Carlo Rovelli.
\newblock Relational quantum mechanics.
\newblock \emph{International Journal of Theoretical Physics}, 35(8):1637--1678, 1996.

\bibitem[Sellars(1956)]{sellars1956}
Wilfrid Sellars.
\newblock Empiricism and the philosophy of mind.
\newblock In Herbert Feigl and Michael Scriven, editors, \emph{The Foundations of Science and the Concepts of Psychology and Psychoanalysis}, Minnesota Studies in the Philosophy of Science, volume~1, pages 253--329. University of Minnesota Press, Minneapolis, 1956.

\bibitem[Spinoza(1677)]{spinoza1677}
Benedict de Spinoza.
\newblock \emph{Ethica, ordine geometrico demonstrata}.
\newblock In \emph{B.d.S. Opera Posthuma}. [Amsterdam: Jan Rieuwertsz], 1677.

\bibitem[Wolfram(2002)]{wolfram2002}
Stephen Wolfram.
\newblock \emph{A New Kind of Science}.
\newblock Wolfram Media, Champaign, IL, 2002.
\end{thebibliography}

\end{document}
